\subsection{Chamber assignment under environment constraints}
\label{subsec:chamber_assignment}

For a fixed sample vector \(\mathbf{s}\), this stage sets each
chamber to one temperature--humidity environment.
Tests are assigned later by the scheduler.
Pull-down work scales with \(\mathbf{s}\), so the assignment is
recomputed whenever the sample vector changes.

The workload \(W_e(\mathbf{s})\) of the problem statement is the
total processing time of every required test that needs
environment \(e\).
Voltage-sensitive work is the same sum, taken only over the
products that require a voltage-capable chamber,
\(\mathcal{P}^v=\{p\in\mathcal{P}:\nu_p=1\}\):
\begin{equation}
\label{eq:workload_voltage}
W_e^V(\mathbf{s})
=
\sum_{p\in\mathcal{P}^v}
\sum_{\substack{t\in\mathcal{T}_p\\e(t)=e}}
m_{pt}(\mathbf{s})\tau_t.
\end{equation}
Only environments with \(W_e(\mathbf{s})>0\) enter the model.
Chamber \(c\) has \(M_c\) stations; \(M^{tot}\) and
\(M^{V,tot}\) are the totals of all stations and of
voltage-capable stations.
The two loads are turned into station targets by sharing each
pool in proportion to its work:
\begin{equation}
\label{eq:targets}
T_e
=
\frac{W_e(\mathbf{s})}{\sum_{e'}W_{e'}(\mathbf{s})}\,M^{tot},
\qquad
T_e^V
=
\frac{W_e^V(\mathbf{s})}{\sum_{e'}W_{e'}^V(\mathbf{s})}\,M^{V,tot}.
\end{equation}
If \(W_e^V(\mathbf{s})=0\), we set \(T_e^V=0\).

The decision variables are the assignment
\(x_{ce}\in\{0,1\}\), equal to \(1\) when chamber \(c\) is set to
environment \(e\); the resulting station counts
\(S_e=\sum_{c\in\mathcal{C}}M_cx_{ce}\) and
\(V_e=\sum_{c\in\mathcal{C}}M_cv_cx_{ce}\); and the deviations
\(z_e,z_e^V\ge 0\).
Each chamber takes one environment, every demanded environment
receives at least one chamber, every environment with
\(W_e^V(\mathbf{s})>0\) receives at least one voltage-capable
chamber, and a humidity setting is forbidden when \(h_c=0\).
The objective is the total deviation from the two targets:
\begin{align}
\min\quad
&
\sum_{e\in\mathcal{E}}z_e+\sum_{e\in\mathcal{E}}z_e^V
\label{eq:assign_obj}\\
\text{s.t.}\quad
&
\sum_{e\in\mathcal{E}}x_{ce}=1
&& \forall c\in\mathcal{C},
\label{eq:assign_one}\\
&
\sum_{c\in\mathcal{C}}x_{ce}\ge 1
&& \forall e\in\mathcal{E},
\label{eq:assign_cover}\\
&
\sum_{c\in\mathcal{C}}v_cx_{ce}\ge 1
&& \forall e\in\mathcal{E}:W_e^V(\mathbf{s})>0,
\label{eq:assign_volt}\\
&
x_{ce}=0
&& \text{if }e\text{ needs humidity and }h_c=0,
\label{eq:assign_hum}\\
&
z_e\ge |S_e-T_e|,
\quad
z_e^V\ge |V_e-T_e^V|
&& \forall e\in\mathcal{E},
\label{eq:assign_dev}\\
&
x_{ce}\in\{0,1\},\quad z_e,z_e^V\ge 0.
&&
\label{eq:assign_dom}
\end{align}

A short local search then adjusts \(\mathbf{x}\).
A swap exchanges the settings of a busy chamber and an idle
chamber; a move copies a busy setting onto an idle chamber.
The first feasible neighbor that reduces the EDD tardiness of
\(\mathbf{s}\) is accepted.
Algorithm~\ref{alg:chamber_assignment} is the procedure.

\begin{algorithm}[H]
\caption{ChamberAssignment}
\label{alg:chamber_assignment}
\begin{algorithmic}[1]
\State \textbf{Input:} sample vector \(\mathbf{s}\)
\State \textbf{Output:} chamber assignment \(\mathbf{x}\)
\State Compute \(W_e(\mathbf{s})\), \(W_e^V(\mathbf{s})\), \(T_e\), and \(T_e^V\) \label{step:ca-load}
\State \(\mathbf{x}\leftarrow\) the optimal solution of \eqref{eq:assign_obj}--\eqref{eq:assign_dom} \label{step:ca-exact}
\Repeat \label{step:ca-loop}
    \State Try load-guided swaps and moves that stay feasible \label{step:ca-nbhd}
    \If{a neighbor improves the EDD tardiness of \(\mathbf{s}\)} \label{step:ca-accept}
        \State Replace \(\mathbf{x}\) by that neighbor
    \EndIf
\Until{no improving feasible neighbor remains} \label{step:ca-stop}
\State \Return \(\mathbf{x}\)
\end{algorithmic}
\end{algorithm}

Step~\ref{step:ca-load} forms the loads and the targets.
Step~\ref{step:ca-exact} solves the assignment model.
Steps~\ref{step:ca-loop}--\ref{step:ca-accept} accept the first
feasible neighbor that improves tardiness, and stop when a pass
finds none (Step~\ref{step:ca-stop}).
